\chapter*{مقدمه}\markboth{مقدمه}{مقدمه}
\addcontentsline{toc}{chapter}{مقدمه}

هدف پروژه‌ی کارشناسی حاضر، مطالعه‌ی دقیق الگوریتم
\lr{Greedy Local K-Means} یا به اختصار \lr{GLKM}
\cite{yuan2026} و بررسی امکان بهبود آن است. این الگوریتم با بازنویسی تابع
هدف \lr{k-means} به‌صورت مجموع فاصله‌های زوجی و با استفاده از یک ماتریس
فاصله‌ی تنک، از ادغام حریصانه‌ی خوشه‌ها استفاده می‌کند. در نقطه‌ی فعلی
پروژه هنوز یک تغییر عمومی و دارای برتری برای \lr{GLKM} به دست نیامده
است.
فعالیت‌های انجام‌شده در سه محور قرار می‌گیرند. محور نخست، مطالعه‌ی منابع
تاریخی و الگوریتمی مهم شامل کارهای مک‌کوئین، لوید، هارتیگان--وونگ و
\lr{k-means++} است. محور دوم،
پیاده‌سازی اجرای \lr{GLKM} برای بازه‌ای از مقادیر $k$ و تولید نمودار eblow
است. محور سوم، پیاده‌سازی ایده‌ی استفاده از افراز اولیه‌ی \lr{k-means++} برای
کاهش تعداد خوشه‌های آغازین \lr{GLKM} است.

برای ملموس‌تر شدن تفاوت الگوریتم‌ها، دو صفحه‌ی وب مستقل و تعاملی نیز به کمک هوش مصنوعی ChatGPT ساخته
شد: یکی مقایسه‌ی گام‌به‌گام روش مک‌کوئین با روش لوید، و دیگری
مقایسه‌ی \lr{Hartigan--Wong} با لوید. مسیر این خروجی‌ها در پوشه‌ی
\lr{interactive} پروژه ثبت شده است.

ساختار گزارش چنین است: فصل نخست مبانی اولیه را توضیح
می‌دهد؛ فصل دوم به روش‌های پس از آن می‌پردازد؛ فصل سوم \lr{GLKM} را تشریح
می‌کند و فصل‌های چهارم و پنجم دو آزمایش و نتایج آن‌ها را
گزارش می‌کنند.
